%% kapa: exam=1 round=36 year=2025 period=1 form=A subject=civil_law subject_ko=민법 seq=27 number=27
%% source: Q-Net 감정평가사 기출문제 (cst00309) artlSeq=5231525 file=raw/1008/5231525/2025년 제36회 감정평가사 1차시험 1교시 문제지 원본.pdf
%% conversion: kordoc 4.24.51 -> tools/gen1.py (auto); review=proofread-sonnet; flags=none
\begin{question}{27}
점유의 회수에 관한 설명으로 옳은 것은? (다툼이 있으면 판례에 따름)
\begin{choices}
  \item 사기에 의한 의사표시에 의해 물건을 인도해 준 경우 점유의 침탈이 인정된다.
  \item 직접점유자가 임의로 점유를 이전한 경우, 그 점유이전이 간접점유자의 의사에 반하면 점유의 침탈이 인정된다.
  \item 점유회수청구권은 침탈자의 악의의 특별승계인에게는 주장하지 못한다.
  \item 간접점유자가 점유회수청구권을 행사하는 경우, 먼저 자기에게 반환할 것을 청구해야 한다.
  \item 점유회수청구권을 행사할 수 있는 기간은 그 기간 내에 소를 제기해야 하는 기간으로 해석된다.
\end{choices}
\end{question}
\answer{5}
